\section{Benchmark and Grounded Corpus}
\label{sec:benchcorpus}
OptSemBench-C is a behavior-contract benchmark: it generates probes that cover optimizer-relevant feature interactions and measures which public contracts those probes activate.  It deliberately does not rank engines by latency.  This choice follows the benchmark principle that the metric and denominator must match the object of study.  CardEst asks whether cardinality estimates improve DBMS plan choice~\cite{han2021cardest}; CardBench stresses learned cardinality estimators in relational settings~\cite{chronis2024cardbench}; Learned Cardinalities models correlated joins~\cite{kipf2019learned}; NeuroCard models probabilistic cardinalities~\cite{yang2019neurocard}; DeepDB learns database distributions~\cite{hilprecht2020deepdb}; FactorJoin specializes join-cardinality estimation~\cite{wu2023factorjoin}; Neo learns plan choices~\cite{marcus2019neo}; and Bao evaluates learned hinting behavior~\cite{marcus2021bao}.  An optimizer-contract benchmark must instead expose semantic traps in the comparison vocabulary.

The contract schema is fixed before source extraction, witness counting, and repair search.  Its fields come from prior optimizer architecture rather than from the measured collisions.  Operator/action and layer follow System R~\cite{selinger1979access}, Volcano~\cite{graefe1990volcano}, Cascades~\cite{graefe1995cascades}, Calcite~\cite{begoli2018calcite}, and Orca~\cite{soliman2014orca}.  Placement follows Garlic federation~\cite{carey1995garlic}, BigDAWG polystores~\cite{stonebraker2015bigdawg}, Dremel serving trees~\cite{melnik2010dremel}, and Snowflake cloud execution~\cite{dageville2016snowflake}.  Decision time follows progressive optimization~\cite{markl2004progressive}, mid-query reoptimization~\cite{kabra1998reoptimization}, adaptive query processing~\cite{deshpande2007adaptive}, eddies~\cite{avnur2000eddies}, and Spark adaptive execution~\cite{armbrust2015spark}.  Observability is admitted from plan, profile, and stage evidence in the public source ledger, with provenance-style evidence separation~\cite{green2007provenance}.  Layer and placement enter the predeclared field universe; exhaustive retained-field testing later reports them as a sufficient headline repair basis, not as fields selected after seeing failures.

Let $F=F_1\times\cdots\times F_n$ be the finite feature domain and $\Psi$ the validity constraints.  A target interaction is a valid partial assignment over a feature subset.  The schema and feature domain are fixed before witness counting and derive from standard query-processing constructs and cited public optimizer surfaces, not from witnesses or motifs.  The benchmark targets global pairwise coverage plus selected three-way interactions for source boundary, connector capability, operator class, statistics need, join shape, adaptivity trigger, and distribution trigger.  The generator builds one deterministic candidate pool from renderable feature interactions and admitted public-rule guards.  It first includes 99 guard-forced completions so that every admitted rule has a reachable query shape, then greedily covers the remaining feature-interaction universe.  These two roles are reported separately: guard-forced probes prevent public contracts from becoming unreachable paper claims, while the remaining probes complete the declared feature-interaction denominator.  Neither role is treated as an independent held-out population sample.  A probe is useful when it activates a contract guard or separates two contract signatures, not when it estimates workload frequency.

\begin{observation}[Coverage]
If every valid target interaction has at least one renderable full assignment, the generator terminates with a probe set covering every renderable target interaction.
\end{observation}
\begin{proof}
The target universe is finite.  Each selected assignment covers at least one uncovered interaction.  Thus the number of uncovered interactions decreases monotonically until it reaches zero.
\end{proof}

The main empirical corpus contains only 287 source-linked grounded rules.  Each rule records a public evidence span, source URL, line range, digest, and source class.  Claims without grounded evidence, expressible guards, or public observability are outside the finite denominator rather than negative evidence.  Thus the paper claims coverage of the declared public-contract denominator, not a census of private optimizer rules.  Figure~\ref{fig:benchmark-snapshot} gives the denominator view before any collision measurement: sources, rules, probes, and motif requirements are fixed, and the zero-failure checks verify support, consistency, SQL execution, and motif coverage.  The corpus covers plan observability, pushdown, join search, distribution, materialization, runtime adaptivity, statistics, and optimizer control surfaces.  Every grounded rule is triggered by at least one generated probe, and contract merge produces no conflicts.

\begin{figure}[t]
\centering
\begin{tikzpicture}[
  stage/.style={vpanel, draw=vldbBlue!66, fill=vldbMist, text width=1.50cm,
    minimum height=.72cm, font=\rmfamily\fontsize{6.45}{7.05}\selectfont,
    text=black!70, align=center, inner xsep=1pt, inner ysep=1pt,
    text height=1.48ex, text depth=.25ex},
  check/.style={vpanel, draw=vldbTeal!74!black, fill=vldbMint,
    minimum width=1.36cm, minimum height=.42cm,
    font=\rmfamily\bfseries\fontsize{6.5}{7.1}\selectfont, text=vldbInk,
    align=center, inner xsep=1pt, inner ysep=.8pt,
    text height=1.42ex, text depth=.22ex},
  axis/.style={vpanel, draw=vldbRule, fill=white, minimum width=1.22cm,
    minimum height=.32cm, font=\rmfamily\fontsize{6.2}{6.8}\selectfont,
    text=vldbGray, align=center, inner xsep=.8pt, inner ysep=.5pt,
    text height=1.35ex, text depth=.20ex},
  zero/.style={vpanel, draw=vldbGold!75!black, fill=vldbWarm,
    minimum width=1.22cm, minimum height=.36cm,
    font=\rmfamily\fontsize{6.25}{6.85}\selectfont, text=black!70,
    align=center, inner xsep=1pt, inner ysep=.8pt,
    text height=1.40ex, text depth=.22ex}
]
\path[use as bounding box] (-.06,-4.10) rectangle (7.25,.34);

\node[vsmall, anchor=west, text=vldbGray] at (.02,.13)
  {denominator construction, not a latency benchmark};
\node[stage] (ev) at (.86,-.56) {\textbf{287}\\evidence\\spans};
\node[stage] (rules) at (2.70,-.56) {\textbf{287}\\admitted\\rules};
\node[stage] (sql) at (4.54,-.56) {\textbf{4,216}\\executable\\SQL probes};
\node[stage] (motif) at (6.38,-.56) {\textbf{90/90}\\published\\motifs};
\draw[vflow] (ev) -- (rules);
\draw[vflow] (rules) -- (sql);
\draw[vflow] (sql) -- (motif);

\node[vsmall, anchor=west, text=vldbGray] at (.02,-1.48)
  {admission checks};
\node[check] at (.86,-1.95) {support};
\node[check] at (2.70,-1.95) {guards};
\node[check] at (4.54,-1.95) {SQL};
\node[check] at (6.38,-1.95) {motifs};
\node[zero] at (.86,-2.44) {0 invalid};
\node[zero] at (2.70,-2.44) {0 conflict};
\node[zero] at (4.54,-2.44) {0 shape};
\node[zero] at (6.38,-2.44) {0 exec};

\node[vsmall, anchor=west, text=vldbGray] at (.02,-3.02)
  {covered query-feature axes};
\node[axis] at (1.02,-3.32) {source};
\node[axis] at (2.30,-3.32) {connector};
\node[axis] at (3.58,-3.32) {join};
\node[axis] at (4.86,-3.32) {predicate};
\node[axis] at (6.14,-3.32) {adapt};
\node[axis] at (1.02,-3.72) {agg.};
\node[axis] at (2.30,-3.72) {stats};
\node[axis] at (3.58,-3.72) {reuse};
\node[axis] at (4.86,-3.72) {dist.};
\node[axis] at (6.14,-3.72) {control};
\end{tikzpicture}
\caption{Corpus, benchmark, and execution denominators. The same grounded denominator carries source evidence, generated probes, admission checks, and published-motif coverage.}
\Description{A compact evidence panel with four denominator counts, four admission checks, observed zero audit or execution issues, and representative covered query-feature axes.}
\label{fig:benchmark-snapshot}
\end{figure}


\begin{table*}[t]
\centering
\caption{Source-to-rule traceability for representative contracts. Each row binds a rule state/action to a public source span, digest prefix, and paraphrase used by the replay manifest.}
\label{tab:evidence-examples}
\footnotesize
\setlength{\tabcolsep}{3pt}
\begin{tabular}{@{}llp{0.15\textwidth}p{0.18\textwidth}p{0.42\textwidth}@{}}
\toprule
\multicolumn{3}{@{}l}{Contract rule} & \multicolumn{2}{l@{}}{Auditable evidence} \\
\cmidrule(r){1-3}\cmidrule(l){4-5}
Engine & State & Action & Span/digest & Source-supported paraphrase \\
\midrule
PostgreSQL~\cite{postgresql_docs} & MAY & Plan/observe/local & PG-EXPLAIN-001 bfb47e66 L46--L56 & EXPLAIN exposes estimated startup and total cost, output rows, row width, and planner-cost units. \\
BigQuery~\cite{bigquery_docs} & MAY & Exchange/adapt/cloud & BIGQUERY-PLAN-001 cc7b206a L990--L992 & The query plan may be modified during execution by adding repartition stages. \\
Snowflake~\cite{snowflake_docs} & MAY & Plan/observe/cloud & SNOWFLAKE-EXPLAIN-001 93a61359 L150--L154 & EXPLAIN compiles a statement and returns a logical execution plan rather than executing it. \\
ClickHouse~\cite{clickhouse_docs} & MAY & Join/adapt/local & CLICKHOUSE-JOINING-001 a1f757e2 L313 & With automatic join-algorithm selection, execution can switch from hash join to partial merge join. \\
DuckDB~\cite{duckdb_docs} & MAY & Materialize/inline/local & DUCK-WITH-001 2a60c8ab L482--L487 & CTE inlining is guarded by single use, absence of volatility, no grouped aggregation, and inlining hints. \\
Spark SQL~\cite{spark_docs} & MAY & Plan/adapt/distributed & SPARK-PERF-001 7e5e2b00 L129--L133 & Adaptive query execution uses runtime statistics and can re-optimize a query while it is running. \\
Trino~\cite{trino_pushdown_docs} & MUST\_NOT & Agg./pushdown/source & TRINO-PUSH-001 cc177219 L168--L176 & Aggregation pushdown does not support aggregate calls containing expressions. \\
\bottomrule
\end{tabular}
\Description{Examples of source-to-rule traceability across seven engines, including state, action, source identifier, digest prefix, line range, and paraphrased evidence.}
\end{table*}


The source set is fixed before witness counting.  We admit an engine only when public sources provide: a comparable plan or contract surface, enough detail to extract an action and applicability guard, and stable analytical-SQL behavior rather than a private vendor note.  PostgreSQL, Trino, BigQuery, Spark SQL, Snowflake, DuckDB, and ClickHouse enter through versioned source records in the evidence ledger; Calcite and DataFusion source records provide schema context rather than denominator engines.  These records are not scholarly support and not performance measurements.  They are the measured public contracts that the framework formalizes.  Proprietary workload logs, undocumented cost-model traces, future engine versions, and operator families outside the declared feature domain are excluded rather than silently counted.

The source-linked denominator spans seven engines: PostgreSQL contributes 80 admitted rules, Trino 61, BigQuery 37, Spark SQL 36, ClickHouse and Snowflake 26 each, and DuckDB 21; the largest single source supports 46 rules.  This design separates admission, probing, and execution validation.  Public sources admit rules; generated probes test guard reachability and SQL executability; DuckDB and PostgreSQL runs validate the probe denominator on two open engines, not every vendor contract.  All-engine execution would conflate SQL validity with public-contract admission.  A DuckDB source span can admit a DuckDB public contract, but a later DuckDB execution is only SQL execution and visible-plan sanity evidence; it never admits rules or chooses a repair field.  Conversely, engines without local execution still participate in collision measurements only through source-linked public contracts, source leave-one-out checks, and side-balanced witness support.

The grounded corpus is intentionally conservative.  A rule is admitted only when the evidence span supports an optimizer action, a guard or applicability condition, and an observability or placement interpretation.  Vague performance advice is excluded from the main corpus.  This policy reduces corpus size but strengthens the central claim: the evaluation is about contracts that are publicly supported, not about inferred implementation behavior.

\textbf{Contract extraction and validation checks.}  Each admitted rule records the engine, version scope, source identifier, digest, action fields, guard, and observability surface.  Human coding maps public evidence spans into the fixed schema; automated checks validate schema consistency, source linkage, guard reachability, trigger coverage, merge conflicts, SQL executability, and table-source consistency.  This is a single-extractor source-line protocol rather than formal verification or recoding agreement: every accepted row is inspectable, vague performance advice is excluded by rule, and excluded rows are outside the declared denominator.  Table~\ref{tab:evidence-examples} gives representative spans; the replay materials carry the complete rule-to-source manifest, source digests, claim-evidence graph, and regenerated table manifests.  The field schema and projection portfolio are replay inputs fixed before collision counting.  Layer plus placement is reported as an interpretable repair basis for headline witnesses after enumerating retained-field vocabularies; the frontier tables give scoped safety and do not claim global safety.  Public source records admit claims; EXPLAIN, profile, or stage surfaces bind exposed plan evidence when available; SQL execution validates the probe denominator rather than private optimizer truth.  A counted witness requires two source-supported reference signatures, projection equality under the declared vocabulary, and separation by the reported repair field.

The denominator dashboard in Table~\ref{tab:denominator-dashboard} checks that this conservatism does not create unreachable contracts or syntactic probes.  A finite-disjunction guard semantics evaluates each admitted rule against the generated probe denominator.  All 287 grounded rules have at least one triggering probe; 252 depend on query features rather than global engine labels; and no guard uses a feature dimension or value outside the benchmark domain.

\begin{table*}[t]
\centering
\caption{Denominator and validation dashboard. Bold marks complete coverage or zero observed SQL-shape/execution failures, not contract firing. The motif-cover prefix is underlined because it is the smallest human-inspection slice.}
\label{tab:denominator-dashboard}
\footnotesize
\setlength{\tabcolsep}{4pt}
\begin{tabular}{@{}llrll@{}}
\toprule
Layer & Check & Value & Scope & Interpretation \\
\midrule
Corpus & grounded rules & \textbf{287} & 26 sources & finite public-contract relation \\
Corpus & probe-supported rules & \textbf{287/287} & median support 102 & every admitted rule is triggerable \\
Corpus & non-global guards & 252 & 3 containments & query-conditional, audited overlaps \\
Guard schema & invalid dimensions & \textbf{0} & fixed feature domain & no out-of-domain guard values \\
SQL corpus & generated probes & \textbf{4,216} & digest \texttt{df107f018e37} & full executable denominator \\
SQL corpus & motif-cover prefix & \underline{48} & human inspection & compact replay slice, not cherry-picking \\
Execution & planned/executed probes & \textbf{4,216/4,216} & 91 distinct plans & every generated probe plans and runs \\
Execution & shape/exec failures & \textbf{0/0} & 18,448 rows & validation checks find no SQL-shape or runtime failures \\
Published motifs & covered requirements & \textbf{90/90} & 12 families & every declared motif has executable probe support \\
\bottomrule
\end{tabular}
\Description{A combined denominator dashboard covering grounded rules, guard checks, SQL probes, execution validation, and published motif coverage.}
\end{table*}


Coverage is reported over two denominators.  The first is benchmark coverage: every targeted renderable valid feature interaction is covered by at least one generated probe.  The second is contract triggering: every grounded rule is applicable to at least one probe.  Together these denominators make the empirical scope explicit.  A rule that cannot be triggered by any probe would indicate a benchmark gap; a feature interaction that cannot be rendered is excluded from the coverage denominator and reported as a modeling constraint.

The benchmark also exposes executable SQL shapes.  Each generated probe has a normalized SQL skeleton and shape checks for joins, filters, grouping, ordering, limits, and reuse structure.  The full 4,216-probe corpus is accompanied by a compact 48-probe motif-cover prefix.  The real-engine motif execution set uses the deterministic union of this motif cover and residual feature-axis maxima; it is not a hand-picked success set, and it is not a replacement for the full evaluation.

A common failure mode in generated benchmarks is to produce feature-complete but non-executable query text.  OptSemBench-C therefore validates every SQL skeleton on a deterministic relational catalog with local, remote, and cross-source namespaces.  This validation is not a performance baseline and is not exhaustive contract-trigger proof.  Its role is stricter and more basic: every generated probe must be parseable, plannable, executable, and shape-consistent with its feature vector.  Table~\ref{tab:denominator-dashboard} reports the validation result: all 4,216 generated probes obtain a plan and execute, with zero SQL-shape failures.  Repeating the full corpus across three deterministic catalog populations yields 12,648 successful probe-catalog executions, and DuckDB/PostgreSQL checks add 8,432 full-corpus engine--probe executions plus 142 motif-representative executions with zero execution failures.

Contract validation is layered.  Guard reachability says that each public rule has a generated probe satisfying its admitted guard; SQL execution says that the denominator is not syntactic fiction.  Visible-plan sanity checks on DuckDB and PostgreSQL expose expected plan tokens at mean rates 0.705 and 0.796 over the full corpus, and 0.756 and 0.819 on motif representatives.  These checks do not reverse-engineer private optimizers, prove every documented behavior fires on every cost-model instance, or cross-validate a source by running the same engine.  A counted witness still requires source-supported reference signatures, guard reachability, projection equality, and repair separation; runtime plan satisfaction is a separate SQL-denominator validation layer.

The same fixed probes support two budgeted evaluation orders.  The cover order is the generator's feature-interaction order.  The diagnostic order greedily maximizes newly triggered grounded rules without adding any new probe; it is a rule-coverage schedule over admitted contracts, not an unseen workload sample.  Table~\ref{tab:benchmark-efficiency} reports both.  The cover order triggers all grounded rules within 100 probes; the diagnostic order triggers 286 of 287 rules within 50 probes and all rules within 100.  Random subsets of 100 probes trigger 81\% of rules on average and 85\% at the 95th percentile.  Thus OptSemBench-C is not merely large enough to cover the surface; it also contains compact diagnostic prefixes that expose almost all grounded contracts.

\begin{table}[t]
\centering
\caption{Probe budgets for triggering grounded contract rules; random columns summarize subset baselines, with brackets giving the 5th--95th percentile interval. Bold marks complete trigger coverage.}
\label{tab:benchmark-efficiency}
\footnotesize
\begin{tabular}{rrrrr}
\toprule
Budget & Cover & Diagnostic & Random mean & Random [5,95] \\
\midrule
50 & .672 & .997 & .707 & [.627,.767] \\
100 & \textbf{1.000} & \textbf{1.000} & .811 & [.770,.850] \\
250 & \textbf{1.000} & \textbf{1.000} & .890 & [.864,.913] \\
500 & \textbf{1.000} & \textbf{1.000} & .920 & [.899,.941] \\
1000 & \textbf{1.000} & \textbf{1.000} & .951 & [.934,.969] \\
2000 & \textbf{1.000} & \textbf{1.000} & .973 & [.958,.986] \\
4216 & \textbf{1.000} & \textbf{1.000} & \textbf{1.000} & \textbf{[1.000,1.000]} \\
\bottomrule
\end{tabular}
\end{table}


A benchmark for public optimizer contracts must connect to existing database workloads without becoming another latency benchmark.  We therefore crosswalk published optimizer and benchmark motifs to OptSemBench-C feature requirements: JOB join-order sensitivity~\cite{leis2015good}, POLAR-style join structure~\cite{justen2024polar}, TPC decision-support benchmarking~\cite{poess2000tpc}, TPC-DS design~\cite{nambiar2006tpcds}, SSB~\cite{oneil2009ssb}, LDBC~\cite{erling2015ldbc}, Wisconsin~\cite{bitton1983wisconsin}, warehouse surveys~\cite{darmont2007benchmarking}, YCSB~\cite{cooper2010ycsb}, and LakeBench~\cite{deng2024lakebench}.  Figure~\ref{fig:external-motifs} reports the full executable denominator: 90 declared requirements across 12 motif families, represented by 71 distinct generated SQL probes because several queries satisfy more than one optimizer-decision requirement.  Five representative anchors illustrate the span: JOB join-order search (8/8 requirements), exact-cardinality testing (9/9), CardEst-style plan-quality sensitivity (8/8), adaptive replanning (9/9), and cross-source joinability (7/7).  Table~\ref{tab:motif-difficulty} gives the 12-family difficulty distribution.  The selected probes are feature-interaction representatives, not replacements for JOB or TPC-H queries; the crosswalk validates optimizer-decision coverage rather than workload runtime behavior.  It remains a frozen audit over optimizer-decision surfaces, not copied benchmark SQL or a latency claim; motifs outside the declared feature domain would be denominator misses.  Because motifs are checked after the schema is fixed and before repair fields are used, this is a non-circular denominator check: a counted collision still requires two source-supported signatures, projection equality, side-balanced support, and repair separation.  We claim contract-denominator adequacy and motif representation, not workload reproduction.

\begin{figure}[t]
\centering
\includegraphics[width=\columnwidth]{generated_figures/fig_external_motifs.pdf}
\caption{Published optimizer-motif denominator from the frozen table of 12 motif families and 90 requirements. A matching-probe span is the minimum--maximum number of OptSem-C probes that trigger a family; teal ranges and diamond medians keep sparse motifs visible.}
\Description{A Python-rendered log-scale range plot showing motif-family coverage and matching-probe spans for twelve optimizer and benchmark families.}
\label{fig:external-motifs}
\end{figure}

\begin{table}[t]
\centering
\caption{External motif difficulty. Matching probes are generated OptSemBench-C probes satisfying each benchmark-family motif requirement.}
\label{tab:motif-difficulty}
\scriptsize
\begin{tabular}{lrrrr}
\toprule
Suite & Motifs & Min & Median & Max  \\
\midrule
TPC-H & 8 & 1 & 87.0 & 1384  \\
TPC-DS & 8 & 86 & 164.5 & 287  \\
JOB & 8 & 55 & 87.0 & 236  \\
SSB & 6 & 89 & 101.0 & 3442  \\
CardEst & 7 & 72 & 228.0 & 259  \\
Adaptive & 6 & 72 & 205.5 & 235  \\
Federated & 8 & 72 & 91.5 & 633  \\
Controls & 8 & 70 & 72.0 & 73  \\
PostgreSQL & 8 & 72 & 73.0 & 73  \\
DuckDB & 8 & 68 & 72.0 & 364  \\
ClickHouse & 7 & 70 & 70.0 & 70  \\
Spark/Trino & 8 & 69 & 70.0 & 236  \\
\bottomrule
\end{tabular}
\Description{Motif difficulty table showing how many probes cover each external benchmark suite motif.}
\end{table}


The motif-difficulty view prevents a misleading denominator.  Some motifs are broad and match many probes, such as star-schema group-by shapes.  Others are sparse, such as a TPC-H conjunctive-filter motif that has one matching rendered probe under the current validity constraints.  Reporting both coverage and matching-probe counts makes benchmark adequacy inspectable: coverage says the motif is present, while difficulty says how concentrated or redundant that coverage is.  The executable motif suite adds a second denominator check: for each published motif, the system selects one satisfying probe and validates the generated SQL on the deterministic catalog, yielding 90 validated motif requirements, 71 distinct representatives, and 383 result rows.  Thus the motif denominator is not a list of names; it is a set of optimizer-decision requirements with concrete SQL representatives and an explicit many-to-one map.  Engine-plan validation is a separate experiment because the motif suite itself is a denominator check rather than a claim that published benchmark queries and generated probes induce identical plans.
